% META: {"id":"TSPE-GEOESPACE-ME-008","chapitre":"TSPE-GEOMETRIE-ESPACE","type_objet":"methode","capacites_codes":["C12"],"methodes":["M8"],"sources_inspiration":[],"status":"approved","origin":"generated","status_history":["generated","audited","accepted","approved"],"release_acceptance":"RELEASE_OWNER_BATCH_ACCEPTANCE","acceptance_closure_digest":"sha256:9b3ccf9a81c5520fbb7e03b7b2d3e7bf2057a02834b6d908bf3be5f49f3b553f","acceptance_receipt_digest":"sha256:4fad86f0282e0fa4e0419cca1ad6379ae7af5a280c04a4a90880f90a1c044242"}
% BEGIN-VERIFY
% from sympy import Matrix
% A=Matrix([1,0,3]); u=Matrix([2,1,-1])
% assert A + 2*u == Matrix([5,2,1])             # N appartient a d
% assert A + 2*u != Matrix([5,2,0])             # M n'y est pas
% END-VERIFY
\begin{fichemethode}{M8}{Déterminer et exploiter une représentation paramétrique}

\quandUtiliser{%
  Pour écrire la représentation paramétrique d'une droite définie par deux
  points ou par un point et une direction, pour tester l'appartenance d'un
  point, et pour reconnaître que deux représentations décrivent la même droite.
}

\pasApas{%
  \begin{enumerate}
    \item \textbf{Identifier} un point $A(x_A;y_A;z_A)$ et un vecteur
          directeur $\vec{u}(\alpha;\beta;\gamma)$, obtenu par différence si
          deux points sont donnés.
    \item \textbf{Écrire} $x=x_A+\alpha t$, $y=y_A+\beta t$,
          $z=z_A+\gamma t$, avec $t\in\mathbb{R}$.
    \item \textbf{Pour tester si $M$ appartient à la droite}~: résoudre la
          première équation en $t$, puis reporter cette valeur dans les
          \emph{deux} autres.
    \item \textbf{Pour comparer deux droites}~: vérifier que les vecteurs
          directeurs sont colinéaires, puis qu'un point de l'une appartient à
          l'autre.
  \end{enumerate}
}

\verifier{%
  Reprendre la représentation obtenue avec deux valeurs de $t$ et contrôler
  que les points trouvés sont bien sur la droite de départ~: pour $t=0$ on
  doit retrouver $A$.
}

\pieges{%
  Une droite admet une infinité de représentations paramétriques~: changer de
  point ou multiplier le vecteur directeur en donne d'autres, tout aussi
  correctes. Et pour l'appartenance, les trois équations partagent le
  \emph{même} paramètre~: un $t$ qui convient pour deux coordonnées sur trois
  ne prouve rien.
}

\exempleRedige{%
  Droite $d$ passant par $A(1;0;3)$ et dirigée par $\vec{u}(2;1;-1)$~:
  \[
    x=1+2t,\qquad y=t,\qquad z=3-t .
  \]
  Le point $N(5;2;1)$ appartient-il à $d$~? De $5=1+2t$ on tire $t=2$~; alors
  $y=2$ et $z=3-2=1$~: les trois équations sont vérifiées avec la même valeur,
  donc $N\in d$. En revanche $M(5;2;0)$ n'y est pas~: il faudrait $z=1$.
}

\end{fichemethode}
